\documentclass[12pt,toc=bib,toc=listof]{scrartcl}

\usepackage[english]{babel}
\usepackage[utf8]{inputenc}
\usepackage[T1]{fontenc}
\usepackage{lmodern}
\usepackage{setspace}
\usepackage{float}
\usepackage{booktabs}
\usepackage{array}
\usepackage{tabularx}
\usepackage{geometry}
\usepackage{caption}
\captionsetup{font=small}
\usepackage{amsmath}
\usepackage{xurl}
\usepackage{microtype}
\usepackage{hyperref}
\usepackage{listings}
\usepackage{graphicx}
\usepackage{ragged2e}
\usepackage{tikz}
\usetikzlibrary{arrows.meta,positioning}
\usepackage{etoolbox}
\AtBeginEnvironment{table}{\singlespacing}
\AtBeginEnvironment{figure}{\singlespacing}

\hypersetup{
  colorlinks=true,
  linkcolor=blue,
  citecolor=blue,
  filecolor=blue,
  urlcolor=blue
}

\lstset{
  basicstyle=\ttfamily\small,
  breaklines=true,
  columns=fullflexible,
  frame=single,
  showstringspaces=false
}

\newcolumntype{P}[1]{>{\RaggedRight\arraybackslash}p{#1}}
\newcolumntype{Y}{>{\RaggedRight\arraybackslash}X}

\newcommand{\hhnsubject}{Deep Learning}
\newcommand{\hhnsubjectnum}{285652}
\newcommand{\hhnlecturer}{Carsten Lanquillon}
\newcommand{\reprttopic}{ClaimGuard: An Evidence-Backed Pipeline for Academic Citation Integrity Checking}

\newcommand{\reprtstudentname}{Wagner Nils}
\newcommand{\reprtstudentid}{212405}
\urldef{\reprtstudentmail}\url{nwagner1@stud.hs-heilbronn.de}

\newcommand{\reprtstudentnameB}{Wenzel Nick}
\newcommand{\reprtstudentidB}{207678}
\urldef{\reprtstudentmailB}\url{nwenzel@stud.hs-heilbronn.de}

\usepackage[headsepline]{scrlayer-scrpage}
\pagestyle{scrheadings}
\clearscrheadfoot
\ihead{\hhnsubject: ClaimGuard}
\ohead{\pagemark}

\titlehead{%
  \flushright
  \IfFileExists{graphics/hhn_en.png}{%
    \includegraphics[width=0.28\linewidth]{graphics/hhn_en.png}%
  }{}%
}
\subject{{\hhnsubject{} (\hhnsubjectnum{})}}
\title{\reprttopic}
\author{
  \reprtstudentname\footnote{\reprtstudentid, \reprtstudentmail} \and
  \reprtstudentnameB\footnote{\reprtstudentidB, \reprtstudentmailB}
}
\publishers{Submitted to \hhnlecturer}

\begin{document}

\pagenumbering{Roman}
\selectlanguage{english}
\maketitle

\newgeometry{left=20mm, top=25mm, right=20mm, bottom=25mm}
\newpage
\pagenumbering{arabic}
\onehalfspacing

\section{Abstract}

Academic integrity checking goes beyond copied text: a genuine reference may fail to support its citing sentence. ClaimGuard identifies claim-bearing sentences and missing citations, validates references, retrieves evidence, and assigns five support labels. Its interface preserves the evidence trail and compares verifiers on identical retrieved input. It combines rules, scholarly metadata APIs, lexical or Sentence-BERT retrieval, local Ollama, an OpenAI frontier model, a SciFact LoRA classifier, and optional Fast-DetectGPT scoring. On a 30-case benchmark, lexical verification achieved $0.900$ accuracy and Macro-F1; heuristic, Ollama, OpenAI, and LoRA achieved Macro-F1 values of $0.868$, $0.728$, $0.518$, and $0.172$. On 50 annotated report sentences, claim-type accuracy fell to $0.520$ and citation-needed F1 to $0.267$. Thirty references produced $0.467$ final-status accuracy; DOI resolution was correct in seven cases, while retraction recall was $0.500$. On 90 aligned SciFact pairs, LoRA beat majority ($0.367$ versus $0.167$ Macro-F1) but trailed Ollama ($0.433$). Fast-DetectGPT reached $0.696$ F1 and $0.924$ AUROC on 15 human and 15 GPT-4 PubMed passages, but missed seven generated passages. ClaimGuard is reviewer triage, not an automatic misconduct detector. Limitations include single-annotator data, synthetic evidence, access bias, public-API limits, and model non-determinism.

\textbf{Keywords:} academic integrity, citation verification, retrieval-augmented generation, scientific claim verification, LoRA, human-in-the-loop review

\section{Introduction}

Academic reports increasingly combine conventional references, web-accessible research, and language-model-assisted writing. A bibliography may contain malformed or fabricated records, but even a genuine paper may not establish the claim attached to it. Plagiarism detection addresses textual overlap; it does not determine whether a sentence needs a citation, whether the cited work exists, or whether available evidence supports the sentence. Manually answering these questions requires repeated movement between a report, bibliographic databases, and source texts, which encourages superficial spot checks as documents grow.

The risk is asymmetric. Missing a misleading citation weakens assessment quality, while a false positive may unfairly imply misconduct. ClaimGuard consequently frames its output as decision support. It prioritizes passages for inspection and exposes the sentence, citation marker, parsed metadata, retrieved evidence, verifier, confidence, latency, and rationale. No score is converted into an accusation.

The project investigates four questions: (RQ1) Can lightweight section-aware rules identify claim types and missing citations in authentic PDFs? (RQ2) Can IEEE-like references be recovered and validated despite PDF extraction errors? (RQ3) How well do lexical and semantic retrieval support a five-way source-support judgment? (RQ4) What trade-offs emerge among a transparent heuristic, local Ollama, an OpenAI frontier model, and a SciFact-trained LoRA adapter?

All three core modules were implemented: claim and citation analysis, reference parsing and validation, and retrieval-based support verification. Optional work comprises interchangeable verifier backends, Sentence-Transformer retrieval, a GPU-trained LoRA adapter, Fast-DetectGPT paragraph scoring, and Streamlit comparison. Fast-DetectGPT was evaluated on paired original and GPT-4 PubMed passages; Binoculars exceeded the 8~GB GPU. The contribution is an inspectable connection between citation need, bibliographic identity, evidence, and alternative verifiers.

\section{Related Work}

SciFact introduced expert-written scientific claims with support, refutation, and evidence rationales \cite{wadden2020}. MultiVerS extended scientific verification with weak supervision and full-document context \cite{wadden2022}. Both show why absence of evidence must remain distinct from contradiction. ClaimGuard adds upstream citation-need detection and bibliographic identity checks to this verification setting.

Citation-needed detection is a distinct upstream task. Redi et al. developed a citation-need taxonomy and models for Wikipedia verifiability \cite{redi2019}; their domain and editorial conventions differ from student reports. Sarol et al. paired retrieval with verification to study subtle citation-integrity errors in biomedical publications \cite{sarol2024}. ClaimGuard targets the same claim--evidence relationship but adds local document ingestion and metadata validation.

Retrieval-augmented generation links models to external evidence \cite{lewis2020}, while Sentence-BERT supports efficient semantic retrieval \cite{reimers2019}. GROBID and S2ORC provide document and bibliographic infrastructure \cite{lopez2009,lo2020}; OpenAlex and Crossref expose complementary scholarly metadata \cite{priem2022,crossref}. scite classifies citation contexts in indexed literature \cite{nicholson2021}, whereas ClaimGuard starts from an arbitrary local report and connects its sentence to evidence.

LoRA enables parameter-efficient adaptation \cite{hu2022}. Binoculars uses cross-model perplexity, while Fast-DetectGPT uses conditional probability curvature \cite{hans2024,bao2024}. Paraphrasing and distribution shift undermine detector reliability \cite{sadasivan2023}; broad tool tests likewise caution against treating such scores as misconduct evidence \cite{weberwulff2023}. ClaimGuard therefore treats AI detection as optional risk signaling. Its gap is integration: citation need, bibliographic identity, retrieval, verifier disagreement, and a reviewer-facing evidence trail in one workflow.

\section{System Architecture and Methods}

\subsection{Pipeline}

\begin{figure}[H]
\centering
\begin{tikzpicture}[
  node distance=5mm,
  block/.style={draw=blue!70!black, rounded corners, fill=blue!4,
    align=center, text width=3.35cm, minimum height=1.05cm,
    font=\scriptsize, inner sep=3pt},
  flow/.style={-{Latex[length=2.2mm]}, thick, blue!70!black}
]
\node[block] (input) {\textbf{PDF/TXT Input}};
\node[block, right=of input] (parser) {\textbf{Parser}\\text extraction and sections};
\node[block, right=of parser] (claims) {\textbf{Claim Extraction}\\citation-presence check};
\node[block, right=of claims] (refs) {\textbf{Reference Parser}\\fields and DOI};
\node[block, below=7mm of refs] (apis) {\textbf{Metadata APIs}\\Crossref, S2, OpenAlex};
\node[block, left=of apis] (rag) {\textbf{Evidence Retrieval/RAG}\\lexical or embeddings};
\node[block, left=of rag] (verify) {\textbf{Verifiers}\\heuristic, local, frontier, LoRA};
\node[block, left=of verify] (ui) {\textbf{Streamlit UI}\\integrity report and audit trail};
\draw[flow] (input) -- (parser);
\draw[flow] (parser) -- (claims);
\draw[flow] (claims) -- (refs);
\draw[flow] (refs) -- (apis);
\draw[flow] (apis) -- (rag);
\draw[flow] (rag) -- (verify);
\draw[flow] (verify) -- (ui);
\end{tikzpicture}
\caption{ClaimGuard architecture and data flow. Evidence, confidence, latency, and backend metadata remain attached to the final report.}
\label{fig:architecture}
\end{figure}

\begin{table}[H]
\centering
\scriptsize
\setlength{\tabcolsep}{3pt}
\begin{tabularx}{\textwidth}{@{}P{1.55cm}YYYY@{}}
\toprule
\textbf{Core module} & \textbf{Implemented components} & \textbf{Output} & \textbf{Evaluation evidence} & \textbf{Limitations} \\
\midrule
1: Claim extraction and citation presence & Section-aware sentence types, citation linking, inherited citations, severity rules & Sentence label, cited indices, missing-citation flag/severity & 50 real sentences: accuracy .520, Macro-F1 .437; citation-needed F1 .267 & Single annotator; rules confuse definitions, methods, and interpretations \\
2: Bibliographic validation & IEEE-like parser; identity/DOI matching; Crossref, Semantic Scholar, OpenAlex; retraction flags & Parsed record, candidate, DOI, confidence, validation/problem status & 30 references: parser 1.000; status Acc. .467/F1 .412; DOI 7/7; retraction F1 .667 & Manual small gold set; 26/30 degraded API coverage; repeated S2/OpenAlex HTTP 429; two retraction positives \\
3: Deep claim--source verification/RAG & Chunking, lexical/embedding retrieval, heuristic/Ollama/OpenAI/LoRA interfaces & Evidence chunks, five-way verdict, rationale, backend metadata & 30 controlled cases: Macro-F1 heuristic .868, Ollama .728, OpenAI .518, LoRA .172 & Synthetic short passages; incomplete full text; monetary cost not logged \\
\bottomrule
\end{tabularx}
\caption{Core Module Compliance. Metrics refer only to the listed saved evaluations.}
\label{tab:compliance}
\end{table}

The parser accepts PDF and plain text, cleans extraction artifacts, separates body, references, and supplied evidence, and propagates section headings. Sentence splitting protects initials, abbreviations, decimal values, and IEEE citation markers. The claim module assigns one of five types: factual claim, methodological statement, opinion/interpretation, background/definition, or non-claim. Rules combine empirical verbs, numerical patterns, method cues, hedging, and definitional constructions. A factual sentence outside methods or results is flagged when it lacks a same-sentence or safely inherited citation; severity rises for numerical, causal, universal, or explicit prior-work assertions.

The reference module extracts index, authors, title, venue, year, and DOI and repairs common PDF artifacts. A parsed DOI is normalized and used for DOI-specific Crossref, Semantic Scholar, and OpenAlex lookups before title search; matching scores combine DOI, title, author, year, and venue. DOI resolution alone cannot prove identity: after a benchmark exposed chimeric citations, exact-DOI boosting was constrained by title and year consistency. Crossref retraction relations/subtypes and OpenAlex's \texttt{is\_retracted} field can produce \texttt{retracted\_or\_problematic}. Coverage remains partial because Semantic Scholar supplies no retraction flag, DOI resolution is not a complete registry audit, and public services can throttle queries. In manual runs on 7 July 2026, Semantic Scholar requests without an API key and OpenAlex requests with a configured \texttt{mailto} parameter repeatedly returned HTTP 429, including after waiting several minutes and changing the input file. This records the observed service condition rather than a systematic API-availability benchmark. Candidate metadata and API failures remain visible rather than being silently accepted.

Evidence from abstracts, supplied passages, metadata, or accessible full text is split into overlapping chunks. Lexical retrieval normalizes tokens, removes stop words, lightly stems terms, and scores overlap. Semantic retrieval encodes claims and chunks with \texttt{all-MiniLM-L6-v2} and normalized cosine similarity. Retrieval is restricted to the cited reference when possible, reducing contamination from unrelated works.

\subsection{Verifier backends and interface}

The deterministic verifier uses coverage, negation, directional conflict, and explicit absence-of-evidence cues. Ollama sends the same claim and top evidence to \texttt{qwen3:1.7b} at temperature zero. The OpenAI backend uses \texttt{gpt-5.4-mini}, disables response storage, and requires structured JSON. Both return exactly one of \emph{supported}, \emph{partially supported}, \emph{not supported}, \emph{contradicted}, or \emph{insufficient evidence}, plus confidence and rationale. The shared prompt instructs models to judge only the supplied evidence; its full text and a sanitized response are in Appendix~\ref{app:prompt}.

The local SciFact classifier starts from \texttt{distilbert-base-uncased} with rank-8, alpha-16 LoRA adapters on query/value projections. Training used 1,261 pairs, 450 validation pairs, batch size 16, three epochs, and 30.8 seconds on an RTX 3070 Ti. The adapter contains 740,355 parameters (1.09\% of 67,696,134 loaded parameters) and occupies 2.83 MiB. A seeded balanced subset (30 per native label) compares it directly with temperature-zero \texttt{qwen3:1.7b} prompting and a majority baseline on identical claim--abstract pairs. Its three native labels map to only three of ClaimGuard's five categories.

\subsection{Evaluation design}

Evaluation is separated into authentic-document and controlled evidence. The authentic claim set contains ten manually selected sentences from each of five reports ($N=50$), labeled once for claim type and citation need. Module~2 uses 20 authentic references plus ten controlled perturbations. Controlled Module~1 tests contain 25 patterns; Module~3 uses 30 balanced synthetic claim--evidence pairs. All four verifiers received identical retrieved chunks. LoRA validation uses all 450 SciFact pairs; its aligned comparison uses a reproducible balanced 90-pair subset. Module~4 uses 15 original PubMed answers and their 15 paired GPT-4 generations from the official Fast-DetectGPT repository (commit \texttt{971b052}, seed 42; both members at least 40 words). This set provides provenance labels but no citation markers, so detector--citation correlation is not evaluated.

\section{Results}

\subsection{Quantitative results}

\begin{table}[H]
\centering\small
\begin{tabularx}{\textwidth}{@{}Y r r r r r@{}}
\toprule
Component & $N$ & Acc. & Prec. & Recall & F1 \\
\midrule
\multicolumn{6}{@{}l}{\textit{Synthetic/controlled regression evidence}} \\
Synthetic citation-needed detection & 25 & 1.000 & 1.000 & 1.000 & 1.000 \\
Five-way verification, lexical & 30 & 0.900 & 0.905 & 0.900 & 0.900$^{a}$ \\
Five-way verification, embeddings & 30 & 0.867 & 0.900 & 0.867 & 0.868$^{a}$ \\
SciFact LoRA validation & 450 & 0.409 & -- & -- & 0.353$^{a}$ \\
SciFact majority baseline & 90 & 0.333 & -- & -- & 0.167$^{a}$ \\
SciFact LoRA, aligned subset & 90 & 0.389 & -- & -- & 0.367$^{a}$ \\
SciFact Ollama zero-shot & 90 & 0.533 & -- & -- & 0.433$^{a}$ \\
Module-2 final status & 30 & 0.467 & -- & -- & 0.412$^{a}$ \\
DOI resolution & 7 & 1.000 & 1.000 & 1.000 & 1.000 \\
Retraction flag & 30 & 0.967 & 1.000 & 0.500 & 0.667 \\
Fast-DetectGPT, PubMed/GPT-4 & 30 & 0.767 & 1.000 & 0.533 & 0.696 \\
\addlinespace
\multicolumn{6}{@{}l}{\textit{Authentic-document evidence}} \\
Real claim-type classification & 50 & 0.520 & -- & -- & 0.437$^{a}$ \\
Real factual-claim detection & 50 & -- & 0.455 & 0.833 & 0.588 \\
Real citation-needed detection & 50 & -- & 0.250 & 0.286 & 0.267 \\
Reference single-record parsing & 30 & 1.000 & -- & -- & 1.000$^{b}$ \\
\bottomrule
\multicolumn{6}{@{}l}{\footnotesize $^{a}$Macro-F1; $^{b}$one parsed record per input. Module-2 gold labels were manually assigned.}
\end{tabularx}
\caption{Saved evaluation results, explicitly separated by evidence type.}
\label{tab:results}
\end{table}

The dominant finding is the synthetic-to-real gap. Handcrafted patterns were perfect on the 25-case regression set, yet real-report claim accuracy was 0.520. Definitions were never recovered correctly as a class, and nine of fifteen methodological statements became factual claims. Citation-needed recall was 0.286: five of seven annotated omissions were missed. Rules remain useful for prioritization but do not generalize sufficiently for automatic enforcement.

Fast-DetectGPT achieved .924 AUROC with 1.019~s mean API latency. It produced no false positives on the 15 human passages, but detected only 8/15 GPT-4 passages. Thus its high precision came with limited recall on this small unedited set; it is a triage signal rather than authorship evidence. The benchmark contains no citation markers, so no association between detector score and citation behavior can be inferred.

\begin{table}[H]
\centering
\scriptsize
\begin{tabular}{@{}lrrrrr@{}}
\toprule
\textbf{Gold $\backslash$ predicted} & \textbf{Factual} & \textbf{Method} & \textbf{Opinion} & \textbf{Definition} & \textbf{Non-claim} \\
\midrule
Factual claim & 15 & 0 & 2 & 0 & 1 \\
Methodological statement & 9 & 6 & 0 & 0 & 0 \\
Opinion/interpretation & 4 & 0 & 5 & 0 & 1 \\
Background/definition & 5 & 0 & 0 & 0 & 2 \\
\bottomrule
\end{tabular}
\caption{Confusion matrix for the 50 authentic sentences. No gold examples were labeled non-claim.}
\label{tab:claim_confusion}
\end{table}

\begin{table}[H]
\centering
\scriptsize
\begin{tabularx}{\textwidth}{@{}P{2.25cm}r r r r Y@{}}
\toprule
\textbf{Module-2 task} & $N$ & \textbf{Prec.} & \textbf{Recall} & \textbf{F1/Acc.} & \textbf{Notes} \\
\midrule
Single-record parsing & 30 & -- & -- & 1.000 Acc. & 20 authentic + 10 controlled references \\
Exists/not found & 30 & 1.000 & .538 & .700 F1 & Conservative: 12 false negatives, no false positives \\
Final four-way status & 30 & -- & -- & .412 Macro-F1 & Accuracy .467; partial-match F1 .000 \\
DOI resolution & 7 & 1.000 & 1.000 & 1.000 F1 & Five valid, two invalid; chimeric identity checked separately \\
Retraction flag & 30 & 1.000 & .500 & .667 F1 & Only two positive cases; one was missed \\
\bottomrule
\end{tabularx}
\caption{Gold-labeled Module-2 evaluation. API coverage was degraded for 26/30 cases, so status metrics are service-condition dependent.}
\label{tab:module2_status}
\end{table}

On the balanced support benchmark, lexical verification reached 0.900 Macro-F1 and embeddings 0.868. On the aligned SciFact subset, LoRA improved over majority by .200 Macro-F1 (.367 versus .167) but trailed Ollama zero-shot (.433). LoRA nevertheless recovered CONTRADICT with .360 F1, whereas Ollama never correctly predicted that class (F1 .000), instead mapping 21/30 contradictions to SUPPORT. LoRA's five-label transfer remained weak (.172 Macro-F1) because it cannot emit \emph{partially supported} or \emph{not supported}. Fine-tuning therefore added class-specific signal, but did not produce an overall zero-shot improvement and did not transfer cleanly across taxonomies.

\subsection{Qualitative behavior}

\begin{table}[H]
\centering\footnotesize
\begin{tabularx}{\textwidth}{@{}p{2.25cm} Y p{3.35cm}@{}}
\toprule
Case & Observed behavior & Interpretation \\
\midrule
RQ5 sentence 12 & Flagged an uncited claim that retrieval is a central RAG bottleneck. & Useful high-severity triage. \\
Report-5 references & Recovered five titles, years, and venues despite collapsed spacing. & Parser repair succeeded on authentic PDF text. \\
Contradiction C1 & Matched ``do not eliminate errors'' against ``eliminate all hallucinations.'' & Negation and direction cues worked. \\
RQ1-02 & Labeled a result interpretation as a factual claim. & Mixed rhetorical roles confuse rules. \\
RQ5 title page & Flagged header text as an uncited factual claim. & Layout noise survived extraction. \\
\bottomrule
\end{tabularx}
\caption{Five representative successes and failures.}
\label{tab:qualitative}
\end{table}

\subsection{Local versus frontier comparison}

\begin{table}[H]
\centering
\scriptsize
\setlength{\tabcolsep}{3pt}
\begin{tabularx}{\textwidth}{@{}P{1.8cm}r r r r P{2.5cm}Y@{}}
\toprule
\textbf{Backend} & $N$ & \textbf{Acc.} & \textbf{Macro-F1} & \textbf{Mean latency} & \textbf{Tokens / cost} & \textbf{Typical failure} \\
\midrule
Heuristic & 30 & .867 & .868 & 0 ms & no tokens & partial $\rightarrow$ supported (2) \\
Ollama \texttt{qwen3:1.7b} & 30 & .733 & .728 & 2744.5 ms & 3775 prompt / 1674 generated; cost not logged & partial $\rightarrow$ supported (3) \\
OpenAI \texttt{gpt-5.4-mini} & 30 & .633 & .518 & 1424.1 ms & 6050 input / 1466 output; cost not logged & not supported $\rightarrow$ insufficient (5) \\
SciFact LoRA & 30 & .300 & .172 & 6.0 ms & no tokens & partial $\rightarrow$ insufficient (6) \\
\bottomrule
\end{tabularx}
\caption{Four-backend comparison on the same 30 gold-labeled cases and identical retrieved evidence. Model calls occurred on 27 cases; three no-evidence cases abstained before invocation.}
\label{tab:local_frontier}
\end{table}

Ollama/OpenAI agreement was .567 (17/30). The local model exceeded the frontier model by .210 Macro-F1, while OpenAI was faster on this hardware/network run. OpenAI's main error was collapsing \emph{not supported} into \emph{insufficient evidence}; Ollama over-predicted full support for partial evidence. In this controlled benchmark, we did not observe a quality advantage that would justify routine frontier use over the local model, especially because monetary cost was not logged. This conclusion is benchmark-specific and should not be read as a general model ranking.

\subsection{Metadata-search and internal baselines}

\begin{table}[H]
\centering\footnotesize
\begin{tabularx}{\textwidth}{@{}r Y r p{3.3cm}@{}}
\toprule
\# & Queried Report-5 work (short title) & Score & Crossref top-1 judgment \\
\midrule
1 & Retrieval-Augmented Generation for Knowledge-Intensive NLP Tasks & .593 & Incorrect later QA variant \\
2 & Sentence-BERT & .867 & Correct \\
3 & RAG for Large Language Models: A Survey & .780 & Incorrect graph-RAG variant \\
4 & Lost in the Middle & .779 & Incorrect ``Found in the Middle'' \\
5 & LLMs Distracted by Irrelevant Context & .459 & Incorrect unrelated security work \\
\bottomrule
\end{tabularx}
\caption{Five-case Crossref metadata-search baseline. Scores are fuzzy title similarity.}
\label{tab:baseline}
\end{table}

Crossref's top candidate was exact for one of five records. Two incorrect near-neighbors exceeded the nominal 0.75 fuzzy threshold because title terms and year overlapped. This metadata-search baseline does not evaluate citation need or claim--source support. A checkpoint-capable OpenAI-without-RAG ablation was implemented, but its 30-case run was not completed after a provider timeout and usage-limit rejection; no metric is reported. No completed scite, SemanticCite, or other external-tool prediction file exists; the CSV is only a template. A full existing-tool benchmark therefore remains future work. ClaimGuard is an orchestration layer over scholarly services, not their replacement.

\subsection{Answers to the research questions}

\begin{table}[H]
\centering
\scriptsize
\begin{tabularx}{\textwidth}{@{}P{1cm}Y@{}}
\toprule
\textbf{RQ} & \textbf{Evidence-grounded answer} \\
\midrule
RQ1 & Only partially. On 50 authentic sentences, claim-type accuracy was .520 (Macro-F1 .437) and citation-needed F1 was .267; lightweight rules are triage, not comprehensive detection. \\
RQ2 & Partially. The 30-reference benchmark parsed one record per input and resolved all seven DOI checks, but final-status accuracy was .467 under degraded API coverage; retraction recall was .500 on two positives. \\
RQ3 & On 30 controlled cases, lexical verification reached .900 Macro-F1 and embeddings .868. Because evidence was short and synthetic, this does not establish real-document verification accuracy. \\
RQ4 & On five-label inputs, Macro-F1 was .868 heuristic, .728 Ollama, .518 OpenAI, and .172 transferred LoRA. On 90 native-label SciFact pairs, LoRA (.367) beat majority (.167) but not Ollama zero-shot (.433). \\
\bottomrule
\end{tabularx}
\caption{Direct answers to RQ1--RQ4.}
\label{tab:rq_answers}
\end{table}

\section{Discussion}

ClaimGuard's strongest property is its end-to-end audit trail: a reviewer can move from sentence and citation marker to metadata candidate, retrieved passage, backend, and rationale. The main empirical lesson is nevertheless caution. Strong controlled verification ($.900$ Macro-F1) coexists with weak real-document claim transfer ($.437$ Macro-F1), so synthetic regression scores must not be read as deployment performance.

\begin{table}[H]
\centering
\scriptsize
\setlength{\tabcolsep}{3pt}
\begin{tabularx}{\textwidth}{@{}P{2.25cm}P{3cm}YP{3.2cm}@{}}
\toprule
\textbf{Failure mode} & \textbf{Cause} & \textbf{Example/evidence} & \textbf{Mitigation} \\
\midrule
Layout false positive & PDF headers and title text enter sentence stream & RQ5 title/header flagged as an uncited claim & Layout-aware parsing; exclude front matter \\
Claim-type confusion & Mixed rhetorical roles and narrow rules & 9/15 method statements predicted factual & Learned section-aware model with abstention \\
Metadata/API failure & Near-neighbors and throttled services & Status Acc. .467; degraded coverage in 26/30 cases; repeated S2/OpenAlex HTTP 429 in manual runs & Require identity agreement; cache/backoff; show candidates \\
Access bias & Paywall/non-OA or abstract-only evidence & No source-level rate measured; missing full text can yield insufficient evidence & Separate access failure from support; licensed/OA retrieval \\
Taxonomy/model bias & SciFact has three labels; prompts may avoid contradiction & LoRA transfer F1 .172; Ollama contradiction F1 .000 & Align labels; class-balanced prompts/data \\
\bottomrule
\end{tabularx}
\caption{Observed or explicitly unmeasured failure modes and mitigations.}
\label{tab:failure_modes}
\end{table}

Access bias deserves special emphasis: paywalled or non-open-access works are harder to verify than sources with complete machine-readable text. Consequently, \emph{insufficient evidence} may describe retrieval, throttling, or access limits rather than an unsupported claim, and may systematically favor open-access literature. The repeated HTTP 429 responses observed for both Semantic Scholar and OpenAlex show that changing the submitted file does not necessarily restore metadata coverage because rate limits apply to the service/client request stream rather than to one document. API timeouts and incomplete deposits add similar missingness. The evaluation is small and single-annotator; Module~2 has only seven DOI checks and two retraction positives, while the local/frontier comparison uses synthetic short evidence.

False positives are consequential in academic assessment. Common knowledge, inherited citations, extraction errors, and cautious scientific language can all trigger incorrect flags. AI-text scores are particularly unsuitable as authorship proof \cite{sadasivan2023,weberwulff2023}. ClaimGuard therefore avoids misconduct labels and requires source inspection. Fast-DetectGPT's perfect precision on 15 human passages must not be generalized: it missed 7/15 generated passages, and edited or paraphrased text was not tested. Citation association also remains unmeasured because the provenance-labeled benchmark contains no citation markers; private student reports were not transmitted to the external API for an exploratory analysis. Future work should add double annotation, authentic claim--source labels, broader DOI/retraction cases, an external-tool benchmark, and a consented, preregistered AI-detection test including edited/paraphrased text and citation behavior.

\label{mainend}
\clearpage
\begin{thebibliography}{16}
\small
\bibitem{wadden2020} D. Wadden et al., ``Fact or Fiction: Verifying Scientific Claims,'' \emph{EMNLP}, pp. 7534--7550, 2020. doi: 10.18653/v1/2020.emnlp-main.609.
\bibitem{wadden2022} D. Wadden et al., ``MultiVerS: Improving Scientific Claim Verification with Weak Supervision and Full-Document Context,'' \emph{Findings of NAACL}, 2022. arXiv:2112.01640.
\bibitem{lewis2020} P. Lewis et al., ``Retrieval-Augmented Generation for Knowledge-Intensive NLP Tasks,'' \emph{NeurIPS}, vol. 33, pp. 9459--9474, 2020.
\bibitem{reimers2019} N. Reimers and I. Gurevych, ``Sentence-BERT: Sentence Embeddings using Siamese BERT-Networks,'' \emph{EMNLP-IJCNLP}, pp. 3982--3992, 2019.
\bibitem{lopez2009} P. Lopez, ``GROBID: Combining Automatic Bibliographic Data Recognition and Term Extraction for Scholarship Publications,'' \emph{ECDL}, pp. 473--474, 2009.
\bibitem{lo2020} K. Lo et al., ``S2ORC: The Semantic Scholar Open Research Corpus,'' \emph{ACL}, pp. 4969--4983, 2020.
\bibitem{priem2022} J. Priem, H. Piwowar, and R. Orr, ``OpenAlex: A Fully-Open Index of Scholarly Works, Authors, Venues, Institutions, and Concepts,'' arXiv:2205.01833, 2022.
\bibitem{crossref} Crossref, ``REST API Documentation,'' 2026. \url{https://www.crossref.org/documentation/retrieve-metadata/rest-api/}.
\bibitem{nicholson2021} J. M. Nicholson et al., ``scite: A Smart Citation Index that Displays the Context of Citations and Classifies Their Intent Using Deep Learning,'' \emph{Quantitative Science Studies}, vol. 2, no. 3, pp. 882--898, 2021.
\bibitem{hu2022} E. J. Hu et al., ``LoRA: Low-Rank Adaptation of Large Language Models,'' \emph{ICLR}, 2022. arXiv:2106.09685.
\bibitem{hans2024} A. Hans et al., ``Spotting LLMs With Binoculars: Zero-Shot Detection of Machine-Generated Text,'' \emph{ICML}, PMLR 235, pp. 17519--17537, 2024.
\bibitem{bao2024} G. Bao, Y. Zhao, Z. Teng, L. Yang, and Y. Zhang, ``Fast-DetectGPT: Efficient Zero-Shot Detection of Machine-Generated Text via Conditional Probability Curvature,'' \emph{ICLR}, 2024. arXiv:2310.05130.
\bibitem{redi2019} M. Redi, J. Morgan, B. Fetahu, and D. Taraborelli, ``Citation Needed: A Taxonomy and Algorithmic Assessment of Wikipedia's Verifiability,'' \emph{The Web Conference}, pp. 1567--1578, 2019. doi: 10.1145/3308558.3313618.
\bibitem{sarol2024} M. J. Sarol, S. Ming, S. Radhakrishna, J. Schneider, and H. Kilicoglu, ``Assessing Citation Integrity in Biomedical Publications: Corpus Annotation and NLP Models,'' \emph{Bioinformatics}, vol. 40, no. 7, btae420, 2024. doi: 10.1093/bioinformatics/btae420.
\bibitem{sadasivan2023} V. S. Sadasivan, A. Kumar, S. Balasubramanian, W. Wang, and S. Feizi, ``Can AI-Generated Text be Reliably Detected?'' arXiv:2303.11156, 2023.
\bibitem{weberwulff2023} D. Weber-Wulff et al., ``Testing of Detection Tools for AI-Generated Text,'' \emph{International Journal for Educational Integrity}, vol. 19, art. 26, 2023. doi: 10.1007/s40979-023-00146-z.
\end{thebibliography}

\clearpage
\appendix
\section{Reproducibility and Prompt Details}
\label{app:prompt}

\subsection{Environment and saved outputs}
Experiments used Python 3.14.3, PyTorch 2.11.0 with CUDA 12.8, an NVIDIA RTX 3070 Ti with 8~GB VRAM, and seed 42 for LoRA training. Primary artifacts are:
\begin{itemize}
  \item \path{outputs/evaluation/rq_claim_evaluation.json}
  \item \path{outputs/evaluation/reference_evaluation.json}
  \item \path{outputs/evaluation/evaluation_lexical.json}
  \item \path{outputs/evaluation/evaluation_embeddings.json}
  \item \path{outputs/evaluation/evaluation_lora.json}
  \item \path{outputs/evaluation/model_comparison_full.json}
  \item \path{outputs/evaluation/module2_validation.json}
  \item \path{outputs/evaluation/lora_artifact_summary.json}
  \item \path{outputs/evaluation/scifact_lora_vs_zeroshot.json}
  \item \path{outputs/evaluation/fast_detect_gpt.json}
  \item \path{outputs/evaluation/README.md}
  \item \path{outputs/evaluation/crossref_baseline_rq5.json}
  \item \path{outputs/models/scifact-lora/training_report.json}
\end{itemize}
API identifiers are retained locally; credentials are excluded.

\subsection{Minimal reproduction command}
From the \path{CapStone Project/} repository root, after installing the base dependencies, the following PowerShell command reproduces the deterministic 30-case five-label benchmark reported as ``Five-way verification, lexical'' in Table~\ref{tab:results}:
\begin{lstlisting}
python -m scripts.run_evaluation --benchmark data/benchmark/claimguard_benchmark.csv --output outputs/evaluation/evaluation_lexical.json --verifier heuristic
\end{lstlisting}
The generated JSON contains aggregate and per-label metrics, the confusion matrix, individual predictions, retrieved evidence, and runtime metadata. This minimal run requires neither an external API key nor Ollama; the model-comparison and API-dependent commands are documented in the repository README.

\subsection{Verifier prompt}
\begin{verbatim}
Classify the claim using exactly one status: supported,
partially_supported, not_supported, contradicted, or
insufficient_evidence. Judge only from the supplied evidence.
Return JSON with status, confidence (0..1), and a concise rationale.
\end{verbatim}

\subsection{LoRA and execution details}
The adapter used rank 8, alpha 16, query/value projections, an effective batch size of 16, three epochs, and 740,355 saved trainable parameters (1.09\% of 67,696,134 loaded parameters). Training took 30.8 seconds; \path{adapter_model.safetensors} is 2,965,276 bytes (2.83 MiB). Additional reproduction commands and response samples are documented in the repository README and stored with the project outputs.

\end{document}
